% ============================================================
% Proposed revision of Section 6: Discussion and Future Directions
% Sources: WSC_CKME_040226.pdf (current draft) + note2 + note3
% Status: DRAFT — update after non-Gaussian experiments complete
% ============================================================

\section{Discussion and Future Directions}
\label{sec:discussion}

In this paper, we propose CKME--DCP, a distribution-free framework for
constructing prediction intervals under heteroscedasticity, implemented
as a two-stage procedure.
Stage~1 fits a conditional CDF estimator via CKME; Stage~2 allocates
additional simulation budget to selected input sites and performs split
DCP to calibrate the nonconformity scores.
Empirically, CKME--DCP attains nominal marginal coverage and yields
lower interval scores than DCP--DR across all tested configurations.
Against hetGP, CKME--DCP is competitive under Gaussian noise and
behaves more reliably under heavy-tailed distributions: hetGP tends to
overcover and produce unnecessarily wide intervals when the Gaussian
assumption is violated.
Regarding Stage~2 budget allocation, expanding site coverage by
increasing $n_2$ is more beneficial than increasing per-site
replications $r_2$.

\paragraph{Score homogeneity as the key mechanism.}
The advantage of CKME--DCP over quantile-first methods is not merely
empirical---it has a structural explanation rooted in the \emph{score
homogeneity} property of CDF-based nonconformity scores.
Under a location-scale model $Y = f(x) + \sigma(x)\varepsilon$, the
abs-median score $s(x,y) = |\hat{F}(y \mid x) - 0.5|$ has a score
distribution that depends on $x$ through the effective bandwidth
$h_{\mathrm{eff}}(x) = h/\sigma(x)$.
With a fixed CV-tuned bandwidth $h$, this ratio varies across the
input space, producing heterogeneous score distributions and a mild
U-shaped coverage profile observed in our experiments.
In contrast, an \emph{adaptive} bandwidth $h(x) = c\,\sigma(x)$
makes $h_{\mathrm{eff}}(x) = c$ constant, so the score distribution
becomes $x$-independent and a single calibration threshold provides
uniform conditional coverage.
Y-space scores (used by DCP--DR) carry an irremovable structural gap:
under heteroscedasticity, the total variation distance between the
score distributions at two input locations $x, x'$ satisfies
$d_{\mathrm{TV}}(F_{s|x}, F_{s|x'}) > 0$ even as $n \to \infty$,
which explains the persistent spatial non-uniformity seen in DCP--DR results.

\paragraph{Group conditional coverage.}
Beyond marginal validity, the empirical coverage of CKME--DCP remains
spatially uniform across input-space bins.
In our experiments ($K = 10$ equal-width bins of $[0, 2\pi]$,
$M = 11$ macroreplications), the maximum bin-level deviation from the
90\% target is at most 0.10 for CKME--DCP across all six simulator
configurations---below the binomial sampling-noise floor of
approximately 0.17 for bin size $n_k \approx 100$, indicating the
structural gap is negligible.
By contrast, DCP--DR exhibits maximum bin deviations up to 0.33,
which exceed the binomial floor by a factor of approximately two and
are attributable to the non-vanishing structural gap described above.
hetGP is spatially flat but globally miscalibrated
(up to 97.5\% on the Gaussian-mixture simulator), reflecting
distributional misspecification rather than spatial non-uniformity.

\paragraph{Future directions.}
Several extensions follow naturally.
%
First, the current experiments use fixed CV-tuned bandwidth $h$.
Implementing a location-adaptive bandwidth $h(x) = c\,\hat\sigma(x)$,
where $\hat\sigma(x)$ is a Stage~1 variance estimate, is expected to
tighten the score homogeneity and improve coverage uniformity,
particularly at the domain boundaries where $\sigma(x)$ is large.
%
Second, the present evaluation considers Gaussian and Student-$t$
noise. A broader study covering skewed (Gamma) and bimodal (Gaussian
mixture) noise families---with both mild and severe departures from
Gaussianity---will characterize how distributional shape affects CDF
estimation and conformal calibration.
%
Third, the relative contributions of the two-stage design's components
(the $S^0$ uncertainty score, the staged data split, the smooth
indicator, and CV hyperparameter tuning) remain to be quantified
individually; a systematic ablation study is planned.
%
Fourth, the current theoretical analysis addresses marginal validity
(Theorem~2) and consistency of the CDF estimator (Theorem~1).
Formalizing the adaptive-bandwidth score homogeneity argument into a
finite-sample conditional coverage bound, and establishing the
corresponding group-conditional guarantee, are directions for future
theoretical work.
%
Fifth, alternative calibration schemes such as Jackknife$+$ and
cross-conformal variants may provide further efficiency improvements
for CKME--DCP, particularly in small-sample regimes.

\end{document}
% ============================================================
